% AI-translated from the original-language LaTeX: printed 252--260 and following unnumbered title leaf [261].
% Principal PDF 259--268; Google PDF 263--272.
\begingroup
\clearpage
\noindent252\hfill\textsc{Correspondence.}\hfill\mbox{}\par\medskip
\noindent how this is to be done. In the present case the aim is to modify the condition equation
\[
\alpha\mathrm d10+\beta\mathrm d12+\gamma\mathrm d13+\delta\mathrm d20+\varepsilon\mathrm d21+\zeta\mathrm d23+\eta\mathrm d30+\theta\mathrm d31+\varkappa\mathrm d32=\lambda
\]
so that it contains all the corrections. This is very easily done. One need only add
\[
\begin{aligned}
x(\mathrm d01-\mathrm d10+\mathrm d12-\mathrm d21+\mathrm d20-\mathrm d02)&=0\\
y(\mathrm d01-\mathrm d10+\mathrm d13-\mathrm d31+\mathrm d30-\mathrm d03)&=0\\
z(\mathrm d12-\mathrm d21+\mathrm d23-\mathrm d32+\mathrm d31-\mathrm d13)&=0
\end{aligned}
\]
with $x$, $y$, $z$ chosen at will. But we are no better off unless we know how to choose them. I say: choose $x$, $y$, $z$ so that the sum of the 12 squares of the coefficients in the resulting equation is a minimum, thus
\[
(x+y)^2+x^2+y^2+(\alpha-x-y)^2+(\beta+x+z)^2+\text{etc.}=\text{Minimum},
\]
from which you can determine $x$, $y$, $z$ with little effort.

It can be proved (though the proof does require some intricate calculation):

I. That if one carried out this operation four times, namely starting the second time from $\frac{10}{12}.\frac{12}{13}.\frac{13}{10}=1$, then from $\frac{20}{21}.\frac{21}{23}.\frac{23}{20}=1$, and the fourth time from $\frac{30}{31}.\frac{31}{32}.\frac{32}{30}=1$, one obtains 4 final equations which, on closer inspection, are identical to one another, i.e. all 13 terms (including the constant term) are proportional in all of them. The advice [to take] the first expansion based on the largest triangle is intended solely to obtain everything in the largest numbers, which are indeed proportional respectively to the areas of the $\Delta\Delta$ $123$, $023$, $013$, $012$. For safety I usually work out all 4 in any case.

For a configuration of the form \tikz[baseline=-.4ex,x=.027cm,y=.027cm]{\coordinate (a) at (0,0);\coordinate (b) at (1,10);\coordinate (c) at (13,11);\coordinate (d) at (16,0);\draw[line width=.8pt] (a)--(b)--(c)--(d)--cycle;\draw[line width=.3pt] (a)--(c) (b)--(d);} two added together give the same sum as the other two. For a configuration of the form \tikz[baseline=-.4ex,x=.027cm,y=.027cm]{\coordinate (a) at (0,0);\coordinate (b) at (8,11);\coordinate (c) at (20,0);\coordinate (o) at (8,5);\coordinate (t) at (8,8.2);\draw[line width=.8pt] (a)--(b)--(c)--cycle;\draw[line width=.3pt] (a)--(o)--(t) (o)--(c);} one is the sum of the other three. I add all 4, of course taking them in such a way that \textls[100]{absolute} addition takes place.

Thus symmetry is given its full due. There is also a way to shorten the work here, which I pass over since this work is in any case quite negligible compared with the whole operation.

\clearpage
\noindent\mbox{}\hfill\textsc{Adjustment of Simple Figures.}\hfill253\par\medskip
II. But symmetry is only one reason. If there were just one such quadrilateral, the great advantage here would be that adjustment in this form does not disturb at all the adjustment of the 4 (of which only 3 are mutually independent) condition equations [of the 2nd kind] supplied by this quadrilateral. If there are several such quadrilaterals that are separate from one another (have no side in common), as happened repeatedly in my system, the same advantage holds for all of them. Then in each case only the angle sums in the \textls[100]{adjacent} $\Delta\Delta$ are disturbed. But I am convinced that even in a system as intricate as yours, this approach would make the adjustment \textls[100]{very considerably faster}.

I believe you will take pleasure in trying this method on even one of your quadrilaterals. It is entirely analogous for 5-sided figures, 6-sided figures, etc.

I had really almost forgotten this method, so I rediscovered it only while inspecting old calculations from 1826. For the later $\Delta\Delta$, as also for this year's measurements in the Bremen area, I proceed somewhat differently, since they do not warrant so much time. I first adjust just the angle sums precisely. I have become so proficient at this that, for example, I need only one hour for the very intricate Bremen survey. Then I turn to the condition equations of the 3rd kind (side ratios), which I treat in such a way that they are satisfied exactly and those angle sums are \textls[100]{not} disturbed again in the slightest. This does not yield the \textls[100]{absolute minimum} of the sum of squares, but in any case comes close to it and brings everything into perfect agreement. For the Westphalian measurements I believe I modified the method so that I went through the operation several times. This can be arranged so that the more often one repeats it, the closer one comes to the absolute minimum sum of squares. It is, however, quite impossible to describe this procedure briefly.

\clearpage
\noindent254\hfill\textsc{Posthumous Papers.}\hfill\mbox{}\par\medskip
\begin{center}[4.]\\[5pt]On the Adjustment of Angles in a Quadrilateral.\end{center}
Let the 4 points be $a$, $b$, $c$, $d$;
\[
\begin{array}{rlllll}
A&\text{twice}&\text{the area}&\text{of}&\text{triangle}&bcd\\
{[}B&\text{»}&\text{»}&\text{»}&\text{»}&cda\\
C&\text{»}&\text{»}&\text{»}&\text{»}&dab\\
D&\text{»}&\text{»}&\text{»}&\text{»}&abc].
\end{array}
\]
Let the angle $bcd$ be denoted by $Ac$, and similarly for the others.

Set
\[
\begin{aligned}
\log\tfrac{\sin Bc.\sin Cd.\sin Db}{\sin Bd.\sin Cb.\sin Dc}&=\alpha\\[4pt]
\log\tfrac{\sin Cd.\sin Da.\sin Ac}{\sin Ca.\sin Dc.\sin Ad}&=\beta\\[4pt]
\log\tfrac{\sin Da.\sin Ab.\sin Bd}{\sin Db.\sin Ad.\sin Ba}&=\gamma\\[4pt]
\llap{\raisebox{-.5ex}{\text{\tiny.}}\hspace{3.5cm}}\log\tfrac{\sin Ab.\sin Bc.\sin Ca}{\sin Ac.\sin Ba.\sin Cb}&=\delta.
\end{aligned}
\]
One then has [if the angle sums in the triangles have already been adjusted:]
\[
\frac\alpha A=\frac\beta B=\frac\gamma C=\frac\delta D=\mu,
\]
\noindent[and approximately, according to Art. 1, if one sets the direction corrections $C^{01}$ and $C^{10}$, etc., now denoted by $\partial ab$ and $\partial ba$, etc., equal to one another, and at the same time writes $A$ for $-T^0$, $B$ for $T'$, $C$ for $-T''$ and $D$ for $T'''$,]
\[
M\mu=\frac{ab^2}{CD}.\partial ab-\frac{ac^2}{BD}.\partial ac+\frac{ad^2}{BC}.\partial ad+\frac{bc^2}{AD}.\partial bc-\frac{bd^2}{AC}.\partial bd+\frac{cd^2}{AB}.\partial cd;
\]
\[
M=\frac{206265}{k},\qquad k=\text{modulus of Briggsian logarithms}
\]
\noindent\hspace*{.7cm}$\log M=5{,}676\,6408$.
\begin{center}Example.\end{center}
\begin{center}
\begin{tabular}{rcccc}
Quadrilateral:&Deister,&Lichtenberg,&Garssen,&Falkenberg.\\
&$d$&$c$&$b$&$a$
\end{tabular}
\end{center}
[The adjustment of the angle sums of the 4 triangles of the quadrilateral in the plane had given the following values.]

\clearpage
\noindent\mbox{}\hfill\textsc{Adjustment of Simple Figures.}\hfill255\par\medskip
\[
\setlength{\arraycolsep}{2pt}
\begin{array}{l|rrr|l|rrr|l|rrr|l|rrr}
[Ab&47^\circ&18'&46''{,}567&[Ba&34^\circ&33'&59''{,}066&[Ca&57^\circ&33'&17''{,}328&[Da&22^\circ&59'&18''{,}262\\
Ac&66&1&16{,}315&Bc&55&34&13{,}347&Cb&99&14&52{,}202&Db&146&33&38{,}770\\
Ad]&66&39&57{,}118&Bd]&89&51&47{,}587&Cd]&23&11&50{,}470&Dc]&10&27&2{,}968\\ \hline
&180&0&0{,}000&&180&0&0{,}000&&180&0&0{,}000&&180&0&0{,}000
\end{array}
\]
[The approximate values of the logarithms of the sides are:
\[
\begin{array}{l@{\;}r@{\qquad}l@{\;}r}
ab\ .\ .\ .&4{,}44961&bc\ .\ .\ .&4{,}78266\\
ac\ .\ .\ .&4{,}93218&bd\ .\ .\ .&4{,}78052\\
ad\ .\ .\ .&4{,}84854&cd\ .\ .\ .&4{,}68605.
\end{array}
\]
This gives:]
\[
\begin{array}{l@{\;}r@{\qquad}l@{\;}r@{\qquad}l@{\;}r}
\log A=&9{,}42951&\sin Cd\ .\ .\ .&9{,}595\,3854&\sin Ca\ .\ .\ .&9{,}926\,2936\\
\log B=&9{,}53459&\sin Da\ .\ .\ .&9{,}591\,6709&\sin Dc\ .\ .\ .&9{,}258\,6170\\
\log C=&9{,}22445&\sin Ac\ .\ .\ .&9{,}960\,8017&\sin Ad\ .\ .\ .&9{,}962\,9423\\ \cline{4-4}\cline{6-6}
\log D=&8{,}97346&&9{,}147\,8580&&9{,}147\,8529
\end{array}
\]
\noindent\makebox[\linewidth][c]{\hspace{.27\linewidth}$[\beta=0{,}000\,0051]$}\par\medskip
\[
\setlength{\arraycolsep}{3pt}
\begin{array}{l@{}r|@{\quad}l@{}r@{\quad}l@{}r@{\quad}l@{}r}
[\beta.\ .\ .\ .\ .\ &4{,}70757&[ab^2\ .\ .\ .\ &8{,}89922&[ac^2\ .\ .\ .\ .\ &\rlap{\raisebox{1.5ex}{\text{\tiny.}}}9{,}86436&[ad^2\ .\ .\ .\ .\ &9{,}69708\\
B\ .\ .\ .\ .\ &\underline{9{,}53459}&CD\ .\ .\ .\ &8{,}19791&BD\ .\ .\ .\ &8{,}50805&BC\ .\ .\ .\ .\ &8{,}75904\\ \cline{4-4}\cline{6-6}\cline{8-8}
\mu\ .\ .\ .\ .\ &5{,}17298&&0{,}70131&&1{,}35631&&0{,}93804\\[5pt]
M\ .\ .\ .\ .\ &\underline{5{,}67664}&bc^2.\ .\ .\ .\ &9{,}56532&bd^2\ .\ .\ .\ .\ &9{,}56104&cd^2\ .\ .\ .\ .\ &9{,}37210\\ 
M\mu.\ .\ ]&0{,}84962&AD\ .\ .\ ]&8{,}40297&AC\ .\ .\ ]&8{,}65396&AB\ .\ .\ .\ ]&8{,}96410\\ \cline{4-4}\cline{6-6}\cline{8-8}
&&&1{,}16235&&0{,}90708&&0{,}40800
\end{array}
\]
\[
\begin{aligned}
7''{,}073={}&5{,}027\partial ab-22{,}715\partial ac+8{,}670\partial ad+14{,}533\partial bc\\
&\hspace{5.5cm}-8{,}074\partial bd+2{,}559\partial cd.
\end{aligned}
\]
[From this one obtains the corrections
\[
\partial ab=5{,}027\varkappa,\qquad\partial ac=-22{,}715\varkappa,\quad\text{etc.,}
\]
\noindent where
\[
899{,}35\varkappa=7{,}073,
\]
\noindent hence
\[
\log\varkappa=7{,}89567-10
\]
\noindent holds.]

\clearpage
\noindent256\hfill\textsc{Posthumous Papers.}\hfill\mbox{}\par\medskip
\[
\begin{array}{r@{\;}c@{\;}r}
\partial ab&=&+0''{,}0395\\
\partial ac&=&-0{,}1786\\
\partial ad&=&+0{,}0682.\\
\partial bc&=&+0{,}1143\\
\partial bd&=&-0{,}0635\\
\partial cd&=&+0{,}0201.
\end{array}
\]
[The angle corrections are therefore
\[
\begin{aligned}
\partial Ab&=-\partial bc+\partial bd=-0''{,}178\\
\partial Ac&=-\partial cd+\partial bc=+0{,}094\\
\partial Ad&=-\partial bd+\partial cd=+0{,}084
\end{aligned}
\]
\begin{center}etc.\end{center}
Taking the given value of the side $cd$ (Lichtenberg---Deister) as the starting value for computing the sides, one thus obtains the following table:]
\begin{center}
\renewcommand{\arraystretch}{1.08}
\begin{tabular}{|l|r@{\quad}r@{\quad}r|r|r|}
\hline
\multicolumn{1}{|c|}{[ Station}&\multicolumn{3}{c|}{Angles}&\multicolumn{1}{c|}{$\log\sin$}&\begin{tabular}{c}Log. of\\the sides\end{tabular} ]\\ \hline\hline
Garssen&$47^\circ$&$18'$&$46''{,}389$&$9{,}866\,3270$&$4{,}686\,0451$\\
Lichtenberg&66&1&$16{,}409$&$9{,}960\,8018$&$4{,}780\,5199$\\
Deister&66&39&$57{,}202$&$9{,}962\,9424$&$4{,}782\,6605$\\[9pt]
Falkenberg&34&33&$59{,}313$&$9{,}753\,8603$&$4{,}686\,0451$\\
Lichtenberg&55&34&$13{,}148$&$9{,}916\,3595$&$4{,}848\,5443$\\
Deister&89&51&$47{,}539$&$9{,}999\,9988$&$4{,}932\,1836$\\[9pt]
Falkenberg&57&33&$17{,}357$&$9{,}926\,2937$&$4{,}780\,5199$\\
Garssen&99&14&$52{,}305$&$9{,}994\,3182$&$4{,}848\,5444$\\
Deister&23&11&$50{,}338$&$9{,}595\,3848$&$4{,}449\,6110$\\[9pt]
Falkenberg&22&59&$18{,}044$&$9{,}591\,6698$&$4{,}782\,6605$\\
Garssen&146&33&$38{,}695$&$9{,}741\,1930$&$4{,}932\,1837$\\
Lichtenberg&10&27&$3{,}261$&$9{,}258\,6204$&$4{,}449\,6111$\\[5pt]
\end{tabular}
\end{center}

\clearpage
\noindent\mbox{}\hfill\textsc{Adjustment of Simple Figures.}\hfill257\rlap{\hspace{.4em}\raisebox{1.3ex}{\tiny.}}\par\medskip
\begin{center}[5.]\hspace{.5em}.\\[5pt]Quadrilateral between 4 Points 1. 2. 3. 4.\end{center}
The condition equation is formed as follows. Let the product of the perpendiculars to side $12$ from points $3$ and $4$ be $p'p'$, etc. Then, denoting the direction of side $12$ by $a$ [its correction by $\mathrm da=\mathrm d12$], setting $p'p'=A$, and proceeding similarly for the others, one has:
\[
0=\frac{\mathrm da}A+\frac{\mathrm db}B+\cdots
\]
\noindent[or]
\[
\begin{aligned}
0={}&\mathrm d12-\tfrac{\sin142\sin321}{\sin314\sin132}\,\mathrm d13+\tfrac{\sin213\sin421}{\sin314\sin142}\,\mathrm d14\\
&+\tfrac{\sin213\sin421}{\sin324\sin132}\,\mathrm d23-\tfrac{\sin214\sin321}{\sin324\sin142}\,\mathrm d24
+\left\{\begin{array}{cl}
&\tfrac{\sin214\sin321}{\sin143\sin432}\,\mathrm d34\\[5pt]
\text{or}&\tfrac{\sin213\sin421}{\sin431\sin243}\,\mathrm d34.
\end{array}\right.
\end{aligned}
\]
\bigskip
\begin{center}[6.]\\[5pt][Adjustment of a Polygon.]\end{center}
If a number $n$ of triangles close around a single point and
\[
\begin{array}{rlllllll}
A&\text{is the correction}&\text{of the sum}&\text{of the angles}&a,a',a''&\text{in the}&\text{first}&\text{triangle}\\
B&\text{»}&\text{»}&\text{»}&b,b',b''&\text{»}&\text{second}&\text{»}\\
C&\text{»}&\text{»}&\text{»}&c,c',c''&\text{»}&\text{third}&\text{»}
\end{array}
\]
\begin{center}etc.,\end{center}
and $S$ is the correction of the sum of the angles $a+b+c+\text{etc.}$ [that lie around the point], one corrects
\[
\begin{array}{r@{\quad}l@{\;}l}
a&\text{by}&\tfrac{n-1}{3n}A-\tfrac1{3n}B-\tfrac1{3n}C-\cdots+\tfrac1nS\\[6pt]
\left.\begin{array}{r}a'\\a''\end{array}\right\}&\text{each by}&\tfrac{2n+1}{6n}A+\tfrac1{6n}B+\tfrac1{6n}C+\cdots-\tfrac1{2n}S
\end{array}
\]
\begin{center}etc.\end{center}
[It is assumed here that one takes into account only the angle equations, and not
\vfill\noindent{\small IX.}\hfill33

\clearpage
\noindent258\hfill\textsc{Posthumous Papers.}\hfill\mbox{}\par\medskip
\noindent the side equation, and that all angle observations have the same weight, 1.]

This rule can also be formulated as follows.

1) First correct each angle in the first triangle by $\frac13A$, in the second by $\frac13B$, etc.

2) Then correct again those that lie around the single point, so that they close.

3) Then also apply half of this second correction, with the opposite sign, to each of the remaining angles.

The mean error still to be expected [for each $a$, $b$, $c$, $\ldots$ is equal to]
\[
m\sqrt{\frac{2n-2}{3n}}
\]
\noindent[and for each $a'$, $b'$, $a''$, $b''$, $\ldots$ it is equal to]
\[
m\sqrt{\frac{4n-1}{6n}}.
\]
[$m$ is the mean error of unit weight.]
\bigskip
\begin{center}[7.]\\[5pt]Weight of Height Determinations.\end{center}
An initial point $P$ is connected to a series of other [points] $p$, $p'$, $p''$, $p'''$, $\ldots$ $p^{(n)}$, and so on, in such a way that the height differences $Pp$, $Pp'$, $pp'$, $pp''$, $p'p''$, $p'p'''$, $p''p'''$, $p''p^{\mathrm{IV}}$, and so on, that is, the differences between each height and the heights of the two preceding and the two following points, have been observed with equal reliability. Then the weight of the determination of the height difference between $P$ and $p^{(n)}$ by the method of least squares is
\[
\setlength{\arraycolsep}{7pt}
\begin{array}{r@{\;}cccccccccc}
n=&0&1&2&3&4&5&6&7&8&9\\[4pt]
\text{Weight:}&\frac11&\frac32&\frac88&\frac{21}{24}&\frac{55}{75}&\frac{144}{224}&\frac{377}{663}&\frac{987}{1932}&\frac{2584}{5576}&\frac{6765}{15950}.
\end{array}
\]
The numerators are the coefficients of the series obtained by expanding the fraction
\[
\frac1{1-3x+xx}
\]

\clearpage
\noindent\mbox{}\hfill\textsc{Adjustment of Simple Figures.}\hfill259\par\medskip
\noindent and are therefore given by the formula
\[
\tfrac15(3r^{n+1}-2r^n-2r^{-n}+3r^{-n-1})
\]
\noindent where the following relations hold:
\[
\begin{gathered}
2r=3+\sqrt5\quad\text{or}\quad r=2{,}618\,0340\\
\log r=0{,}417\,9752
\end{gathered}
\]
\noindent Likewise, the denominators arise from expanding the fraction
\[
\tfrac{(1-x)^2}{(1+x)(1-3x+xx)^2}
\]
\noindent and are therefore represented by the formula
\[
\tfrac1{25}\{(3n+7)r^{n+1}-(2n+2)r^n-(2n+2)r^{-n}+(3n+7)r^{-n-1}+(-1)^n8\}.
\]
The limiting expression for the weight becomes
\[
\frac{25}{5n+5+4\sqrt5}.
\]
\begin{center}\rule{3cm}{.4pt}\end{center}
\bigskip
\begin{center}\textsc{Remarks.}\end{center}
\begingroup\footnotesize
Note [1] and most of [2], as well as note [7], are in the same notebook. [6] belongs to another notebook. The final result of [2] is on the last page of \textsc{Gauss}'s copy of \textsc{Carnot}'s Geometrie der Stellung, übersetzt von \textsc{Schumacher}. Erster Theil, 1810.

If in a quadrilateral the side equation is transformed with the aid of 3 angle equations, as indicated under [3] in the letter to \textsc{Gerling} of 19 January 1840, it is easy to show (cf. II, p. 253) that the corrections resulting from the adjustment of the quadrilateral are made up of the corrections required to adjust the 3 angle equations alone and those required to adjust the transformed side equation alone. For if one forms the expressions for the correlates from the angle equations and the transformed side equation and uses these expressions to form the normal equations, the normal equations corresponding to the angle equations are \textls[100]{not} coupled with the normal equation arising from the transformed side equation. This presupposes that the factors by which the angle equations were multiplied before being added to the original side equation are chosen so that the sum of the squares of the coefficients in the resulting side equation is a minimum. The constant terms of the angle equations used to transform the side equation need not have been reduced to zero by adjustment beforehand. But if this has been
\par\endgroup
\vfill\noindent\hfill33\textsuperscript{$\ast$}\hspace*{2cm}

\clearpage
\noindent260\hfill\textsc{Remarks.\quad Adjustment of Simple Figures.}\hfill\mbox{}\par\medskip
\begingroup\footnotesize
\noindent done, the constant term of the side equation must of course be computed with the values yielded by the adjustment of the angle equations.

The note under [4] is taken from a calculation notebook belonging to the Hanover arc measurement. The side equations in the form given were set up there in an adjustment of the triangles from the years 1821 to 1823 as far as the side Hamburg-Timpenberg. The triangle angle sums are adjusted first, followed by the adjustment of the abbreviated condition equations arising from the side ratios. The triangle angle sums are unchanged by this adjustment. The end of the letter given under [3], addressed to \textsc{Gerling} on 19 January 1840, presumably refers to this procedure. For numerical calculations one obtains this side equation just as conveniently by forming the coefficients of its corrections in the usual way, either by means of logarithmic differences or by means of the cotangents of the corresponding angles, and then setting the correction of the direction $hi$ equal to the correction of the direction $ih$.

The first formula of Art. [5] is in the same calculation notebook as [4], while the second formula is entered on the front endpaper of a table of logarithms. If the product of the perpendiculars from points $3$ and $4$ to side $12$ is denoted by $P_{12}$, the product of the perpendiculars from $2$ and $4$ to $13$ by $P_{13}$, etc., and if $\mathrm d12$, $\mathrm d13$, etc., denote the changes in direction of $12$, $13$, etc., then, when the quadrilateral is adjusted, the side equation according to Art. [4] reads in approximate form:
\[
\frac{\mathrm d12}{P_{12}}-\frac{\mathrm d13}{P_{13}}+\frac{\mathrm d14}{P_{14}}+\frac{\mathrm d23}{P_{23}}-\frac{\mathrm d24}{P_{24}}+\frac{\mathrm d34}{P_{34}}=0.
\]
The second formula of Art. [5] follows at once. In the original, the term with $\mathrm d14$ in that formula has the minus sign, while the remaining terms are all positive. However, \textsc{Gauss} has added to this formula: “Subject to the signs.”

The following should also be noted about note [7]. If
\[
\llap{\raisebox{1ex}{\text{\tiny.}}\hspace{4cm}}-N_{i-1}+3N_i-N_{i+1}=0,
\]
\noindent where $N_0=0$ and $N_1=1$, the reciprocal of the weight of the height difference is obtained by the formula
\[
\frac{(5n+17)N_{n+1}-8N_n+(-1)^n8}{25N_{n+1}}
\]
\noindent from which, since $\displaystyle\lim_{n=\infty}\frac{N_n}{N_{n+1}}=\frac{3-\sqrt5}{2}$, the limiting expression
\[
\frac{5n+5+4\sqrt5}{25}
\]
\noindent follows.\hfill\textsc{Krüger.}\par
\endgroup
\vspace{1.6cm}
\begin{center}\rule{4cm}{.4pt}\end{center}

\clearpage
% Unnumbered title leaf following printed 260; no folio or signature is visible.
\null\vspace*{\stretch{4}}
\begin{center}{\LARGE STATION ADJUSTMENTS.}\end{center}
\vspace*{\stretch{6}}
\endgroup
% End of the unnumbered title leaf [261].
